\section{Faculty Enablement}\label{sec:faculty}

Faculty are the primary drivers of effective AI use in courses. This section removes barriers and provides support.

\subsection{Permitted and encouraged uses}
Faculty may use AI tools without prior approval (Tier~A, Section~\ref{sec:tools}) for, among others:
\begin{description}[leftmargin=0.5cm]
  \item[Teaching and instruction] lesson planning; generating examples, case studies, worked solutions, analogies, and practice items; adapting explanations to different levels and languages; building course tutors or study assistants configured on the faculty member's own materials.
  \item[Assessment design] drafting rubrics, question banks, and scenario-based or authentic tasks; designing AI-resilient and AI-integrated assessments; generating formative feedback prompts; piloting assessments against AI to see how they perform.
  \item[Content creation] drafting slides, handouts, modules, quizzes, simulations, code samples, and multimedia; translating or localizing materials; producing accessible formats (captions, alt text, plain-language versions).
  \item[Course management and administration] drafting announcements and correspondence, scheduling, organizing course files, summarizing meetings, preparing reports, and using agents to automate routine, non-sensitive workflow steps.
  \item[Feedback and student support] drafting comments on student work for faculty review and edit, building FAQs, and creating personalized study plans, provided the data-handling rules in Section~\ref{sec:privacy} are followed.
  \item[Research] literature exploration, coding and data-analysis support, drafting and editing, and idea generation, subject to Section~\ref{sec:research}.
\end{description}

\subsection{Faculty responsibilities (kept deliberately light)}
\begin{enumerate}
  \item Review and take responsibility for AI-assisted content before use; check facts, citations, and bias.
  \item Do not enter personally identifiable student information, grades, or confidential records into non-approved tools (Section~\ref{sec:privacy}).
  \item AI may assist in grading or feedback, but a faculty member remains accountable for grades and final decisions. Fully automated high-stakes grading is a Tier~C use.
  \item Declare the AI-use level for each course (Section~\ref{sec:integrity}).
\end{enumerate}

\subsection{Simple disclosure norms for faculty}
Disclosure is not required for routine drafting, editing, or planning support. A one-line note (e.g., ``Prepared with AI assistance and reviewed by the instructor'') is expected when AI-generated material is presented to students as substantive instructional content, is published, or when students interact directly with an AI system set up for the course.

\subsection{Professional development and support}
\begin{itemize}
  \item \textbf{Starter pathway:} a short, self-paced orientation on AI capabilities, limits, and responsible use, plus the quick reference in Appendix~\ref{app:quickref}.
  \item \textbf{Practice community:} faculty learning circles, show-and-tell sessions, and a shared repository of tested prompts, course agents, assessment designs, and case studies.
  \item \textbf{Mentoring and micro-grants:} peer mentors and small grants or recognition (e.g., for teaching innovation) for faculty who pilot and share AI-enabled course designs.
  \item \textbf{Workload and recognition:} AI-related course redesign, training, and scholarship of teaching count toward professional development and recognition processes. \todoverify{Alignment with faculty evaluation and promotion rules.}
  \item \textbf{Well-being:} support to address technostress and overload associated with rapid technological change; participation in training is encouraged, not penalized.
\end{itemize}

\subsection{Tool and agent access guidance}
\begin{itemize}
  \item The IT Office maintains a short, regularly updated guidance list of tools by tier, with notes on what data each may receive.
  \item The University will work toward providing institutional access (e.g., licensed or education-tier assistants and agents) so that access does not depend on individual ability to pay. \todoverify{Budget and procurement feasibility.}
  \item Where institutional access is unavailable, faculty may use personally held accounts for Tier~A purposes only, and are not required to purchase tools.
  \item Faculty must not require students to buy paid tools or to create accounts that expose their personal data (Section~\ref{sec:equity}).
\end{itemize}
